% answers_short.tex —— 1985 年·全国卷·文科 · 简短答案速查
% 由 tools/make_short_answers.py 自 answers.tex 的 \gksolhead/\gkanswer 提取，
% 并与 1985/1985_数学真题_全国文_答案要点.md 逐题比对后生成，勿手改（改源 MD 后重跑）。

\begin{gknotice}
\textbf{使用说明}：本节只给最终答案便于快速核对，完整推导见解析版。答案为编写者独立推导，\textbf{非官方标准答案}。
\end{gknotice}

\gkpart{题目}
\begin{itemize}[leftmargin=2.6em]
  \item[\textbf{1.}] （D）$\dfrac{a^3}{6}$
  \item[\textbf{2.}] （A）必要条件
  \item[\textbf{3.}] （C）$\{1,3,7,8\}$
  \item[\textbf{4.}] （B）$y=|\sin x|$
  \item[\textbf{5.}] （B）$78$ 个
  \item[\textbf{6.}] $[-2,1)\cup(1,2]$
  \item[\textbf{7.}] $e=2$
  \item[\textbf{8.}] 最大值为 $2$，最小值为 $-2$
  \item[\textbf{9.}] $64$
  \item[\textbf{10.}] $-8i$
  \item[\textbf{11.}] $S_n=\dfrac{n(2n^2+1)}{3}$
  \item[\textbf{12.}] 恒等式成立
  \item[\textbf{13.}] $x=0$
  \item[\textbf{14.}] $x\in\left[-\dfrac52,2\right)$
  \item[\textbf{15.}] 底面内切圆半径为 $h\cot\beta$
  \item[\textbf{16.}] 对称圆方程为 $(x-4)^2+(y+2)^2=1$，即 $x^2+y^2-8x+4y+19=0$
  \item[\textbf{17.}] $\displaystyle\lim_{n\to\infty}T_n=\begin{cases}1,&0<q\le1,\\[2pt]\dfrac1q,&q>1.\end{cases}$
\end{itemize}
